\section{Zadanie 7}
\autor{Aleksander Wieczorek}
Udowodnij, że \(A_0\) nie jest zamknęty na branie nieskończonych przecięć.
\subsection*{Rozwiązanie}
Zbiór zamknięty na branie przecięć to taki, że dla każdego podzbioru elementów tego zbioru ich przecięcie również należy do tego zbioru.\\
Żeby dostać tezę wystarczy pokazać, że istnieje nieskończony podzbiór przedziałów otwartych, których przecięcie to przedział, który jest domknięty z którejś strony.\\
Możemy wziąć np.
\[\bigcap_{n=1}^\infty \left(-\frac{1}{n},\frac{1}{n}\right)=\{0\}=[0,0] \notin A_0\]
Zbiór \(A_0\) nie jest zamknięty na przecięcia, bo dostaliśmy element spoza niego.\\
jeśli chcemy wziąć zbiór inny niż singleton możemy wziąć np.
\[\bigcap_{n=1}^\infty \left(-\frac{1}{n},1+\frac{1}{n}\right)=[0,1]  \notin A_0\]